\section{Recovery}\label{sec:recovery}

A ghost's key and chunk map exist only in its record. Losing the data folder, or the computer,
would therefore lose every ghost, although its chunks are still on GitHub. GhostVault offers two
ways to keep the records elsewhere, both encrypted in one format: \emph{recovery files}, written on
request to a place the user chooses, and the \emph{repository backup}, kept up to date automatically
in the repositories if the user turns it on. \tabref{tab:loss} summarises what each kind of loss
costs.

\begin{table}[t]
\caption{What a loss costs}\label{tab:loss}
\centering\footnotesize
\begin{tabular}{@{}L{19mm}L{30mm}L{30mm}@{}}
\toprule
\textbf{Lost} & \textbf{Without a bundle} & \textbf{With a recovery file or the backup} \\
\midrule
one record & that ghost & restored from the bundle \\ \addlinespace[1.6pt]
the data folder or the computer & every ghost & every ghost in the bundle \\ \addlinespace[1.6pt]
the settings & repositories must be added again; ghosts unaffected & the same \\ \addlinespace[1.6pt]
a token & nothing: enter a new one & the same \\ \addlinespace[1.6pt]
a repository & every ghost with a chunk in it & the same: a bundle holds keys, not data \\ \addlinespace[1.6pt]
the passphrase or recovery key & --- & that bundle, for good: there is no reset \\
\bottomrule
\end{tabular}
\end{table}

\subsection{The bundle format}
A bundle (\lstref{lst:bundle}) is a fixed first line, a header of one line of JSON, and the
ciphertext. The payload holds every complete ghost's record, normalised to \state{ready}, and the
repository list without tokens; it holds no previews, log or journal notes. Ghosts whose remote copy
is incomplete or being removed are left out. The key is derived from the secret with Argon2id,
256\,MiB, 3 passes and 4 lanes: the second recommended setting of RFC~9106~\cite{rfc9106} with four
times its memory. Where the \code{cryptography} build has no Argon2id, scrypt with $N=2^{17}$, $r=8$,
$p=1$ is used. The parameters are read from the header and capped (at most 2\,GiB, 20 passes,
16 lanes), so a malicious file cannot make GhostVault allocate unbounded memory, and the whole header
is authenticated as associated data. The plaintext is padded to a multiple of 64\,KiB, so a bundle's
size shows roughly, not exactly, how many ghosts it holds.

\begin{lstlisting}[caption={A recovery bundle; the header is shown wrapped.},label={lst:bundle}]
GHOSTVAULT-RECOVERY\n
{"check":"<8-byte HMAC, base64>","cipher":"AES-256-GCM",
 "format":"ghostvault-recovery","kdf":{"name":"argon2id",
 "iterations":3,"lanes":4,"memoryKiB":262144,
 "salt":"<16 bytes, base64>"},"kind":"file",
 "nonce":"<12 bytes, base64>","version":1}\n
AES-256-GCM(key, nonce,
  length (8 bytes) || payload JSON
    || random padding to a multiple of 64 KiB,
  associated data = the two lines above)
\end{lstlisting}

The \code{check} value is the first eight bytes of HMAC-SHA256 under the derived key. It lets
GhostVault tell a wrong passphrase from a damaged file, and lets an attacker test a guess no faster
than decrypting would. Nothing can open a bundle whose secret is lost.

\subsection{Recovery files}
A recovery file is protected by a passphrase of at least 12 characters, or by one GhostVault
generates. It is written under a temporary name in the chosen folder, flushed, read back and
decrypted to check it, and only then given its name, without replacing an existing file unless the
user confirmed. Importing one validates every record exactly like a local one before anything is
written, so a bundle with one malformed record imports nothing. Repositories not yet configured are
added without a token and marked as needing one. A ghost already present is skipped, and one with the
same id but a different record is reported and not imported, so local records are never replaced.
GhostVault tracks which ghosts are in any recovery file, and Settings shows how many are in none.

\subsection{The repository backup}
Turning the backup on generates a \emph{recovery key}: 26 characters drawn at random from Crockford's
base32 alphabet~\cite{crockford} (130 bits) and two check characters that catch typing errors, shown
in seven groups of four. It is shown once and stored nowhere; the user writes it down. GhostVault
derives the backup key from it once, with Argon2id and a random salt, and keeps the derived key in its
settings, so later backups need no slow derivation. After every change (a ghost added or removed,
renamed, a preview hidden or shown, an import) it writes a new version of the bundle, of kind
\code{repository}, to \code{ghostvault/recovery/backup.gvrb} in every usable repository.

The bookkeeping survives crashes because nothing about it is kept only in memory. Each version's
sequence number is reserved durably before anything is sent. A version is recorded as verified, with
a digest of its content and the ids it holds, only after the branch of each repository it was
written to has been read back and holds exactly the uploaded bytes; a repository that missed it
keeps the backup due. Whether a backup is \emph{due} is computed by comparing that digest with the
current records, so a crash before, during or after a push leaves it due, and the next operation or
start pushes it. A ghost counts as protected by the backup only if it is in the
last verified version. A failed backup is shown in Settings and retried; it never fails the change
that caused it. On a new computer the user adds the repositories with tokens and enters the recovery
key; GhostVault takes the version with the highest sequence number found in any repository.

The backup has a cost. Anyone who can read the repositories holds the ciphertext and, in the
backup, the encrypted keys, so their safety rests entirely on the recovery key. And Git's history
keeps every earlier version: turning the backup off, or restoring a ghost, removes the current file
from the branch, not the old versions. Whoever later obtains the recovery key and the repository's
history can decrypt the keys of ghosts removed long ago, whose chunks are also still in that history.

\subsection{Incomplete records}\label{sec:incomplete}
A new backup replaces the previous one. If it left out a ghost because that ghost's record could not
be read, the good version on GitHub would be replaced by one without that ghost's key: the encrypted
chunks would remain, and no current backup could open them. GhostVault therefore lets a backup leave a
ghost out only if GhostVault itself removed it, and it records the ids of the ghosts it removes before
deleting their records. Then:
\begin{itemize}
\item a record that cannot be read, or a record that a backup or recovery file holds but that
GhostVault never removed, stops every backup and every recovery-file export, and the last verified
backup stays as it is;
\item before a repository's backup is replaced, the current one is read and decrypted; if it holds
ghosts this computer does not know (another computer, an older copy of the data folder), it is not
replaced in any repository, and one this computer cannot open is not replaced either, unless the user
chose to replace it when turning the backup on;
\item the user resolves the situation explicitly: by repairing the records from the backup or a
recovery file (an unreadable record is moved to \code{damaged/}, never deleted), or by setting them
aside, in which case the backup carries their last good records forward.
\end{itemize}
Ghosting and restoring go on working meanwhile; only backups and exports pause, and the window says
why. One case remains open: a repository that cannot be read when a new version is written is not
compared, so if it alone held a newer version made elsewhere, the version written to the other
repositories can supersede it on a new computer.
